\newpage
\thispagestyle{empty}

\begin{par}
    \setlength{\baselineskip}{5mm}

    \begin{center}
        МИНИСТЕРСТВО НАУКИ И ВЫСШЕГО ОБРАЗОВАНИЯ РФ \\
        ФЕДЕРАЛЬНОЕ ГОСУДАРСТВЕННОЕ АВТОНОМНОЕ ОБРАЗОВАТЕЛЬНОЕ УЧРЕЖДЕНИЕ ВЫСШЕГО ОБРАЗОВАНИЯ \\
        <<\MakeUppercase{\universitylongsplit}>> \\
        (\MakeUppercase{\university}, \universityshort) \\
        \faculty

        \leftskip=0.4cm
        Кафедра: $\underset{\text{(название кафедры)}}{\text{\dotuline{\department\hspace{0.59\textwidth}}}}$\dotuline{\hfill} \\
        Направление подготовки: \programnum\thinspace\programname \\
        Направленность (профиль): \thinspace{\specialization}
    \end{center}

    \begin{flushright}
        УТВЕРЖДАЮ \\[3mm]
        Зав. кафедрой $\underset{\text{(фамилия, И., О.)}}{\text{\dotuline{\aplastname\thinspace\apinitials}}}$ \\[3mm]
        $\underset{\text{(подпись)}}{\text{\dotuline{\hspace{25mm}}}}$ \\[1mm]
        \datetemplate
    \end{flushright}

    \begin{center}
        \textbf{ЗАДАНИЕ} \\
        \textbf{НА ВЫПУСКНУЮ КВАЛИФИКАЦИОННУЮ РАБОТУ БАКАЛАВРА}
    \end{center}

    \begin{flushleft}
        \leftskip=0.4cm
        Студент\textit{у(ке)} $\underset{\text{(фамилия, имя, отчество, номер группы)}}{\text{\dotuline{\studentfortask\hspace{0.28\textwidth}}}}$\dotuline{\hfill}, группы \dotuline{\group} \\
        Тема $\underset{\text{(полное название выпускной квалификационной работы)}}{\text{\dotuline{\parbox[t]{0.82\textwidth}{\topic}}}}$ \\
        утверждена распоряжением проректора по учебной работе от \dotuline{\hfill}\thinspace № \dotuline{\hspace{15mm}} \\
        Срок сдачи студентом готовой работы: \dotuline{\hfill}20\dotuline{\hspace{5mm}} г. \\
        Исходные данные (или цель работы):
        {
            Публикации, описывающие нотацию музыкальных произведений, представление биологических данных, методы машинного обучения.\hfill
        } \\[1mm]
        Структурные части работы:
        {
            Изучение литературы. Модификация существующей формальной музыкальной нотации для генерации аудио данных. Анализ полученных результатов и написание текста работы.\hfill
        } \\[1mm]
    \end{flushleft}

    \begin{flushleft}
        \leftskip=0.4cm
        \begin{minipage}[t]{0.35\textwidth}
            Руководитель ВКР \\
            \safullname \\
            \sadegree \\
            \sainfotemplate \\[1mm]
            \datetemplate
        \end{minipage}
        \hspace{25mm}
        \begin{minipage}[t]{0.35\textwidth}
            Задание принял к исполнению \\
            \studentinfotemplate \\[1mm]
            \datetemplate
        \end{minipage}
    \end{flushleft}

\end{par}
